\documentclass{article}
	\usepackage{graphicx}%插入图片
	\usepackage{titlesec}%修改标题格式
	\usepackage{amsmath}%大部分数学符号
	\usepackage{mathrsfs}%花体
	\usepackage{amssymb}%双线体mathbb
	\usepackage{xeCJK}%中文支持
	\usepackage{ctex} %中文字体
	\usepackage{geometry}%几何绘图
	\usepackage{setspace}%设置行距
	\usepackage{multicol}%表格合并
	\usepackage{multirow}%表格合并
	\usepackage{subfigure}%次级图片
	\usepackage{hyperref}%调整标注与引用
	\usepackage{indentfirst}%首行缩进}
\usepackage{tikz}%tikz绘图
\usepackage{tikz-feynman}
\usepackage{listings}%代码插入
\lstset{language=Matlab}
\usepackage{xcolor}
\usepackage{longtable}%大型表格跨页显示
\usepackage{booktabs}
% \usepackage{mathpazo} %palatino equations
\usepackage{pgfplots}
\usepackage{siunitx}
\usepackage{mhchem}
\linespread{1.25}
\begin{document}
\begin{flushright}
    规范场
    
    hw2

    2024310867 王亚朋
\end{flushright}
\begin{enumerate}
    \item \begin{equation}
      \mathcal{L}=\bar{\psi}i\gamma^\mu D_\mu\psi -\frac{1}{2}\mathtt{Tr}(F_{\mu\nu}F^{\mu\nu})+\frac{\xi}{2}(\partial_\mu A^{a\mu})^2-c\bar{c}\partial_\mu D^\mu c.
    \end{equation}
    \begin{itemize}
      \item [a.] for $\delta A_\mu=\theta D_\mu c$, $\delta \psi=ig\theta c \psi$, $\delta c^a=-\frac{1}{2}g\theta C^{abc} c^b c^c$, $\delta \bar{c}=\theta B$, $\delta B=0$, we have
      \begin{equation}
        \begin{aligned}
          \delta\left(D_\mu\psi\right)=&ig\theta\left(\partial_\mu c \psi+cD_\mu\psi\right)+ig\theta(\partial_\mu c+ ig[A_\mu, c])\psi,\\
          \delta\left((\partial_\mu A^{a\mu})^2\right)=&2\theta\partial_\nu D^{a\nu}\partial_\mu A^{\mu}\\                       
          \delta\left(-\bar{c}\partial_\mu D^\mu c\right)=-\left(\delta\bar{c}\right)\partial_\mu D^\mu c-\bar{c}\partial_\mu D^\mu \delta c
        \end{aligned}
      \end{equation}
        可以得到 $\delta\mathcal{L}=0.$
        \item [b.]
        \begin{equation}
          \delta\Phi=\theta s\Phi
        \end{equation}
        则
        \begin{equation}
          \begin{aligned}
            sc=&-\frac{1}{2}gf^{abc}c^bc^c,\\
            s^2c^a=&-\frac{1}{2}gf^{abc}\left(sc^bc^c-c^bsc^c\right)=0.
          \end{aligned}
        \end{equation}
        对于Gauge Field 
        \begin{equation}
          s^2A_\mu=s(D_\mu c)=D_\mu(sc)+g\left[sA_\mu,c\right]=0.
        \end{equation}
        对于费米子
        \begin{equation}
          s^2\psi=ig(sc)\psi-g^2cc\psi=0.
        \end{equation}
        因此可以得到
        \begin{equation}
          s^2\Phi=0.
        \end{equation}
    \end{itemize}
\end{enumerate}
\end{document}
